\newcommand{\etas}{\ensuremath{\eta_{\mathrm{s}}}}

\chapter{Introduction}

\section{Motivation}
 
Modern cryptographic systems increasingly rely on the ability to prove
that a computation was carried out correctly, without revealing the
data involved. This idea sits at the centre of zero-knowledge proofs,
verifiable computation, and blockchain-based privacy systems. One of
the most useful building blocks for all of these is the
\textbf{polynomial commitment scheme}.
 
A polynomial commitment scheme lets a prover commit to a polynomial
$f \in \mathbb{F}[X]$ and later convince a verifier that $f$ evaluates
to a specific value at a point of the verifier's choice, without
exposing the polynomial itself. This primitive appears inside nearly
every practical SNARK system in use today, including PLONK, Marlin,
and Groth16 variants.
 
Even though the theory behind these schemes is well established, their
real-world performance depends heavily on how they are implemented and,
in particular, which elliptic curve is used underneath. Different
curves trade off security level, pairing cost, and field arithmetic
speed in different ways. There is surprisingly little published work
that measures all of these effects together in a controlled setting.
This thesis fills that gap by implementing three polynomial commitment
schemes from scratch and benchmarking all four, including Dory via its
reference implementation, across three elliptic curves and measuring
every major performance dimension.
 
\section{Objectives}
 
This thesis has four concrete goals:
 
\begin{enumerate}
    \item \textbf{Implementation.} Build working implementations of
          KZG, Multilinear KZG (MKZG), and Bulletproofs over the
          elliptic curves BLS12-381, BN254, and BLS12-377. For Dory,
          the industry-standard reference implementation by a16z
          \cite{a16z_dory} is used directly, as it is the most
          well-audited and widely adopted codebase for this scheme.
 
    \item \textbf{Benchmarking.} Measure the following metrics for
          each scheme and curve combination:
          \begin{itemize}
              \item Trusted setup time,
              \item Commitment time,
              \item Proving time,
              \item Verification time,
              \item Scalar multiplication time,
              \item Pairing operation time,
              \item Proof size in bytes.
          \end{itemize}
 
    \item \textbf{MSM Optimisation.} Implement Pippenger's
          multi-scalar multiplication algorithm and measure how much
          faster it is compared to the naive approach.
 
    \item \textbf{Analysis.} Use the benchmark data to draw practical
          conclusions about which scheme and curve combination works
          best in different application scenarios.
\end{enumerate}
 
\section{Contributions}
 
The main contributions of this work are:
 
\begin{itemize}
    \item A unified benchmarking framework that covers KZG, MKZG, and
          Bulletproofs from scratch implementations, together with
          Dory via the a16z reference codebase \cite{a16z_dory},
          enabling fair comparison under identical hardware conditions.
 
    \item A complete set of benchmark results covering every major
          performance metric, collected under identical hardware
          conditions.
 
    \item An empirical study of Pippenger's algorithm applied to the
          MSM sub-routine, which is the dominant cost in most
          pairing-based schemes.
 
    \item Practical guidance on choosing the right scheme and curve
          for real applications, backed by measured data rather than
          asymptotic estimates alone.

    \item All implementations and benchmark drivers produced in this
          thesis are openly available at
          \url{https://github.com/farru2610/MTP-II}.
\end{itemize}
 
\section{Thesis Organisation}

\textbf{Chapter~\ref{chap:background}} introduces the
mathematical foundations: finite fields, elliptic curve arithmetic,
bilinear pairings, and the three curves used throughout this work
(BLS12-381, BN254, BLS12-377). It also defines polynomial commitment
schemes formally and states the security properties each construction
must satisfy.
\\
\textbf{Chapter~\ref{chap:kzg}} presents
the KZG scheme for univariate polynomials. It covers the construction,
pseudocode, Rust implementation, and benchmark results across all three
curves, including the impact of Pippenger's MSM optimisation on commit
and prove time.
\\
\textbf{Chapter~\ref{chap:mkzg}} extends KZG to multilinear polynomials. It describes the
tensor-product SRS construction, the bookkeeping-table witness
algorithm, and benchmarks showing how setup cost grows exponentially
in the number of variables $\mu$.
\\
\textbf{Chapter~\ref{chap:bulletproofs}} describes
the transparent, pairing-free commitment scheme based on inner-product
arguments. The chapter covers the recursive halving construction,
Pippenger optimisation for commit and prove, and the $O(n)$
verification cost that limits its practical range.
\\
\textbf{Chapter~\ref{chap:dory}}
presents the Dory scheme, which achieves $O(\log n)$ verification
without a trusted setup using structured $\mathbb{G}_T$ reductions.
The a16z reference implementation is used and benchmarked across all
three curves.
\\
\textbf{Chapter~\ref{chap:comparison}}
brings all four schemes together. It compares setup, commit, prove,
verify, and proof-size metrics side by side, analyses the effect of
Pippenger's algorithm and curve choice, and provides concrete guidance
for selecting the right scheme and curve for a given application.

